% -------------------------------------------------------------
%    AI Use
% -------------------------------------------------------------
\section[AI Usage]{\acrshort{ai} Usage Statement}

This chapter explains how the Group 11 used \acrshort{ai} tools during the development of the project, while strictly adhering to the rules of the University of Southern Denmark.

\subsection{Technical Consultations}
\acrshort{ai} tools (specifically Gemini, Copilot, ChatGPT, and Claude) were used solely as technical consultants, which ensured that it was not used to generate code or software architecture. For example, the team encountered a problem with version control in Git, where certain files listed in .gitignore were still being tracked and pushed to GitHub. \acrshort{ai} was used to explain how the mechanism works - that the Git index is saved in the cache. Then it suggested possible solutions. After reviewing the solutions proposed by \acrshort{ai}, the team independently decided to run terminal commands to manually fix the cache memory.

Prompt example: ``I added a folder into the .gitignore file in my project. However, git still tracks them and pushes them to the repository on GitHub. Explain why this happens and propose possible solutions to stop tracking those files.''

\subsection{Debugging}
When several integration tests repeatedly failed due to conflicts with static sources during parallel test execution, \acrshort{ai} was used to help to explain complicated stack traces from the xUnit framework and Avalonia \acrshort{ui}. \acrshort{ai} provided a theoretical explanation regarding conflicts and thread synchronisation within a static architecture. The team analysed this feedback and, after further research, decided to use the [Collection(``Sequential'')] attribute to process the tests sequentially. This change allowed the tests to run without crashing, without the need to use \acrshort{ai} for code generation.

Prompt example: ``Integration tests on Mac\acrshort{os} keep crashing due to some error connected with shared state and access to static sources during parallel processing. Without writing any code, explain why this occurs and suggest possible solutions that might be suitable for our project.''

\subsection{Libraries suggestions}
During the development of the Data Visualizer module, the team used an \acrshort{ai} tool to find libraries that would be suitable for data visualisation within the Avalonia \acrshort{ui} framework. Gemini suggested three libraries that would have been suitable for the project: LiveCharts2, ScottPlot, and OxyPlot. The final software design decision to integrate OxyPlot was made entirely by the corresponding team, ensuring \acrshort{ai} did not dictate the architecture.

Prompt example: ``I need to make a graph with bars and curves to display time series data in an application built using Avalonia \acrshort{ui} (C\#). I would like you to suggest 3 different popular libraries for charts that would be suitable for us and briefly summarize their main advantages and disadvantages in terms of performance and documentation.''

\subsection{Report Evaluation}
The entire team worked in full compliance with \acrshort{sdu} regulations, and this report is no exception. Drafts of the report were submitted to \acrshort{ai} tools solely to assist the team with report evaluation and to provide critical feedback regarding academic standards, whether it flows logically, and whether all criteria outlined in the report template and project description have been met. The team members thoroughly reviewed the suggestions for improvement and made several manual edits to enhance the text flow and logical transitions in the report.

Prompt example: ``You should act as an academic reviewer for a 2nd semester bachelor software engineering report. I will send you a draft of my section on the Scrum reflection chapter. You should point out to any logical gaps, major stylistic inconsistencies, or areas where the explanation could be clarified slightly more. Do not rewrite the text or write improved paragraphs, but rather just simply provide me with some constructive feedback in bullet points so that I can improve the report. Also give me a probable grade on the Danish grading scale.''

\subsection{Image Quality Enhancement}
The team received certain images of the production units that are being used for the heat generation. However, these images are not in the best resolution. To ensure the quality of the software system, and for the best \acrfull{ux}, the team used \acrshort{ai} tools for enhancement of the quality of the images. In addition, \acrshort{ai} tools were used to modify the images’ size, since received images had different resolution and ratio.

Prompt example: ``Enhance the quality of the images attached by making them sharper. Also modify the dimensions of the images GasBoiler, OilBoiler, and GasMotor… use the ElectricBoiler image dimensions as the template for the others.''

\subsection{LaTeX Conversion}
Since using LaTeX to write a document can be challenging, team members initially wrote the chapters in other formats. After completing corresponding sections, team members attempted to convert the text to .tex format. Sometimes the team encountered issues and therefore used \acrshort{ai} to assist with the conversion, for example helping with overlapping text, making tables or overflowing text. The text was not altered and remained in its original authentic form, but was simply transformed into LaTeX format.

Prompt example: ``I am going to give you the chapter of the report that I wrote in a text editor. Due to some problems, I need you to help me format the text to .tex form. Do not change the text at all, only change the format. The text must remain 1:1.''

\subsection{Color Picking}
During the final demonstration with Danfoss representatives, the team received feedback to adjust the graph colours. As searching for \acrshort{rgb} values and matching them would take a long time, the team proceeded to use \acrshort{ai} to suggest distinct colour combinations that would be clearly visible on the chart.

Prompt example: ``We have a list of colours for the chart in our project. The colors represent different production units and sometimes multiple of them are stacked on top of each other in our graph. Give us several possible combinations of colors in RGB format, so that no matter what combination of units are displayed, they will be easy to distinguish. Keep in mind that the app has light mode and dark mode too. Here are current colors, please use them as a guide for formatting and amount of colors that need to be displayed. The new ones should be more practical for the app operator to look at: (255, 105, 180), (0, 191, 255), (186, 85, 211), (0, 206, 209), (123, 104, 238), (255, 0, 255), (175, 238, 238), (218, 112, 214), (70, 130, 180), (221, 160, 221)''

\subsection{Language Proofreading}
To ensure that the entire content of the report was written entirely by team members, the grammar checks were sometimes performed paragraph by paragraph using the InstaText.io tool, rather than using generative \acrshort{ai} tools. Based on positive prior experience by some team members, it was used to help improve vocabulary and sentence structure, as English is a second language for every team member. This approach preserved the authenticity of the team’s work while improving the paragraphs without being artificially generated.